% Test document: ctex with system fonts, bibliography, a sub-file and a figure
\documentclass[UTF8,fontset=mac]{ctexart}
\usepackage{amsmath,booktabs,graphicx,hyperref}
\newcommand{\soh}{state of health}


\title{测试}
\begin{document}
\maketitle
\tableofcontents
\section{第 1 节 电池健康状态估计}\label{sec:1}
锂离子电池的健康状态估计是电池管理系统的核心功能，电池的容量衰减与内阻增长直接影响电动汽车的续航里程\cite{a2020}。The state of health of a battery is difficult to measure; the data-driven approach uses the charging curve to estimate it, see Eq.~\eqref{eq:1} and Section~\ref{sec:1}.

锂离子电池的健康状态估计是电池管理系统的核心功能，电池的容量衰减与内阻增长直接影响电动汽车的续航里程\cite{a2020}。The state of health of a battery is difficult to measure; the data-driven approach uses the charging curve to estimate it, see Eq.~\eqref{eq:1} and Section~\ref{sec:1}.

锂离子电池的健康状态估计是电池管理系统的核心功能，电池的容量衰减与内阻增长直接影响电动汽车的续航里程\cite{a2020}。The state of health of a battery is difficult to measure; the data-driven approach uses the charging curve to estimate it, see Eq.~\eqref{eq:1} and Section~\ref{sec:1}.

锂离子电池的健康状态估计是电池管理系统的核心功能，电池的容量衰减与内阻增长直接影响电动汽车的续航里程\cite{a2020}。The state of health of a battery is difficult to measure; the data-driven approach uses the charging curve to estimate it, see Eq.~\eqref{eq:1} and Section~\ref{sec:1}.

\begin{equation}\label{eq:1}
  Q_t = Q_0 - \int_0^t \frac{I(\tau)}{3600}\,\mathrm{d}\tau
\end{equation}
\begin{table}[h]\centering\begin{tabular}{lcc}\toprule 方法 & RMSE & MAE\\\midrule GRU & 1.2 & 0.9\\\bottomrule\end{tabular}\caption{结果 1}\end{table}

\section{第 2 节 电池健康状态估计}\label{sec:2}
锂离子电池的健康状态估计是电池管理系统的核心功能，电池的容量衰减与内阻增长直接影响电动汽车的续航里程\cite{a2020}。The state of health of a battery is difficult to measure; the data-driven approach uses the charging curve to estimate it, see Eq.~\eqref{eq:2} and Section~\ref{sec:2}.

锂离子电池的健康状态估计是电池管理系统的核心功能，电池的容量衰减与内阻增长直接影响电动汽车的续航里程\cite{a2020}。The state of health of a battery is difficult to measure; the data-driven approach uses the charging curve to estimate it, see Eq.~\eqref{eq:2} and Section~\ref{sec:2}.

锂离子电池的健康状态估计是电池管理系统的核心功能，电池的容量衰减与内阻增长直接影响电动汽车的续航里程\cite{a2020}。The state of health of a battery is difficult to measure; the data-driven approach uses the charging curve to estimate it, see Eq.~\eqref{eq:2} and Section~\ref{sec:2}.

锂离子电池的健康状态估计是电池管理系统的核心功能，电池的容量衰减与内阻增长直接影响电动汽车的续航里程\cite{a2020}。The state of health of a battery is difficult to measure; the data-driven approach uses the charging curve to estimate it, see Eq.~\eqref{eq:2} and Section~\ref{sec:2}.

\begin{equation}\label{eq:2}
  Q_t = Q_0 - \int_0^t \frac{I(\tau)}{3600}\,\mathrm{d}\tau
\end{equation}
\begin{table}[h]\centering\begin{tabular}{lcc}\toprule 方法 & RMSE & MAE\\\midrule GRU & 1.2 & 0.9\\\bottomrule\end{tabular}\caption{结果 2}\end{table}

\section{第 3 节 改了电池健康状态估计}\label{sec:3}
锂离子电池的健康状态估计是电池管理系统的核心功能，电池的容量衰减与内阻增长直接影响电动汽车的续航里程\cite{a2020}。The state of health of a battery is difficult to measure; the data-driven approach uses the charging curve to estimate it, see Eq.~\eqref{eq:3} and Section~\ref{sec:3}.

锂离子电池的健康状态估计是电池管理系统的核心功能，电池的容量衰减与内阻增长直接影响电动汽车的续航里程\cite{a2020}。The state of health of a battery is difficult to measure; the data-driven approach uses the charging curve to estimate it, see Eq.~\eqref{eq:3} and Section~\ref{sec:3}.

锂离子电池的健康状态估计是电池管理系统的核心功能，电池的容量衰减与内阻增长直接影响电动汽车的续航里程\cite{a2020}。The state of health of a battery is difficult to measure; the data-driven approach uses the charging curve to estimate it, see Eq.~\eqref{eq:3} and Section~\ref{sec:3}.

锂离子电池的健康状态估计是电池管理系统的核心功能，电池的容量衰减与内阻增长直接影响电动汽车的续航里程\cite{a2020}。The state of health of a battery is difficult to measure; the data-driven approach uses the charging curve to estimate it, see Eq.~\eqref{eq:3} and Section~\ref{sec:3}.

\begin{equation}\label{eq:3}
  Q_t = Q_0 - \int_0^t \frac{I(\tau)}{3600}\,\mathrm{d}\tau
\end{equation}
\begin{table}[h]\centering\begin{tabular}{lcc}\toprule 方法 & RMSE & MAE\\\midrule GRU & 1.2 & 0.9\\\bottomrule\end{tabular}\caption{结果 3}\end{table}

\section{第 4 节 电池健康状态估计}\label{sec:4}
锂离子电池的健康状态估计是电池管理系统的核心功能，电池的容量衰减与内阻增长直接影响电动汽车的续航里程\cite{a2020}。The state of health of a battery is difficult to measure; the data-driven approach uses the charging curve to estimate it, see Eq.~\eqref{eq:4} and Section~\ref{sec:4}.

锂离子电池的健康状态估计是电池管理系统的核心功能，电池的容量衰减与内阻增长直接影响电动汽车的续航里程\cite{a2020}。The state of health of a battery is difficult to measure; the data-driven approach uses the charging curve to estimate it, see Eq.~\eqref{eq:4} and Section~\ref{sec:4}.

锂离子电池的健康状态估计是电池管理系统的核心功能，电池的容量衰减与内阻增长直接影响电动汽车的续航里程\cite{a2020}。The state of health of a battery is difficult to measure; the data-driven approach uses the charging curve to estimate it, see Eq.~\eqref{eq:4} and Section~\ref{sec:4}.

锂离子电池的健康状态估计是电池管理系统的核心功能，电池的容量衰减与内阻增长直接影响电动汽车的续航里程\cite{a2020}。The state of health of a battery is difficult to measure; the data-driven approach uses the charging curve to estimate it, see Eq.~\eqref{eq:4} and Section~\ref{sec:4}.

\begin{equation}\label{eq:4}
  Q_t = Q_0 - \int_0^t \frac{I(\tau)}{3600}\,\mathrm{d}\tau
\end{equation}
\begin{table}[h]\centering\begin{tabular}{lcc}\toprule 方法 & RMSE & MAE\\\midrule GRU & 1.2 & 0.9\\\bottomrule\end{tabular}\caption{结果 4}\end{table}

\section{第 5 节 新标题电池健康状态估计}\label{sec:5}
锂离子电池的健康状态估计是电池管理系统的核心功能，电池的容量衰减与内阻增长直接影响电动汽车的续航里程\cite{a2020}。The state of health of a battery is difficult to measure; the data-driven approach uses the charging curve to estimate it, see Eq.~\eqref{eq:5} and Section~\ref{sec:5}.

锂离子电池的健康状态估计是电池管理系统的核心功能，电池的容量衰减与内阻增长直接影响电动汽车的续航里程\cite{a2020}。The state of health of a battery is difficult to measure; the data-driven approach uses the charging curve to estimate it, see Eq.~\eqref{eq:5} and Section~\ref{sec:5}.

锂离子电池的健康状态估计是电池管理系统的核心功能，电池的容量衰减与内阻增长直接影响电动汽车的续航里程\cite{a2020}。The state of health of a battery is difficult to measure; the data-driven approach uses the charging curve to estimate it, see Eq.~\eqref{eq:5} and Section~\ref{sec:5}.

锂离子电池的健康状态估计是电池管理系统的核心功能，电池的容量衰减与内阻增长直接影响电动汽车的续航里程\cite{a2020}。The state of health of a battery is difficult to measure; the data-driven approach uses the charging curve to estimate it, see Eq.~\eqref{eq:5} and Section~\ref{sec:5}.

\begin{equation}\label{eq:5}
  Q_t = Q_0 - \int_0^t \frac{I(\tau)}{3600}\,\mathrm{d}\tau
\end{equation}
\begin{table}[h]\centering\begin{tabular}{lcc}\toprule 方法 & RMSE & MAE\\\midrule GRU & 1.2 & 0.9\\\bottomrule\end{tabular}\caption{结果 5}\end{table}

\section{第 6 节 电池健康状态估计}\label{sec:6}
锂离子电池的健康状态估计是电池管理系统的核心功能，电池的容量衰减与内阻增长直接影响电动汽车的续航里程\cite{a2020}。The state of health of a battery is difficult to measure; the data-driven approach uses the charging curve to estimate it, see Eq.~\eqref{eq:6} and Section~\ref{sec:6}.

锂离子电池的健康状态估计是电池管理系统的核心功能，电池的容量衰减与内阻增长直接影响电动汽车的续航里程\cite{a2020}。The state of health of a battery is difficult to measure; the data-driven approach uses the charging curve to estimate it, see Eq.~\eqref{eq:6} and Section~\ref{sec:6}.

锂离子电池的健康状态估计是电池管理系统的核心功能，电池的容量衰减与内阻增长直接影响电动汽车的续航里程\cite{a2020}。The state of health of a battery is difficult to measure; the data-driven approach uses the charging curve to estimate it, see Eq.~\eqref{eq:6} and Section~\ref{sec:6}.

锂离子电池的健康状态估计是电池管理系统的核心功能，电池的容量衰减与内阻增长直接影响电动汽车的续航里程\cite{a2020}。The state of health of a battery is difficult to measure; the data-driven approach uses the charging curve to estimate it, see Eq.~\eqref{eq:6} and Section~\ref{sec:6}.

\begin{equation}\label{eq:6}
  Q_t = Q_0 - \int_0^t \frac{I(\tau)}{3600}\,\mathrm{d}\tau
\end{equation}
\begin{table}[h]\centering\begin{tabular}{lcc}\toprule 方法 & RMSE & MAE\\\midrule GRU & 1.2 & 0.9\\\bottomrule\end{tabular}\caption{结果 6}\end{table}

\section{第 7 节 电池健康状态估计}\label{sec:7}
锂离子电池的健康状态估计是电池管理系统的核心功能，电池的容量衰减与内阻增长直接影响电动汽车的续航里程\cite{a2020}。The state of health of a battery is difficult to measure; the data-driven approach uses the charging curve to estimate it, see Eq.~\eqref{eq:7} and Section~\ref{sec:7}.

锂离子电池的健康状态估计是电池管理系统的核心功能，电池的容量衰减与内阻增长直接影响电动汽车的续航里程\cite{a2020}。The state of health of a battery is difficult to measure; the data-driven approach uses the charging curve to estimate it, see Eq.~\eqref{eq:7} and Section~\ref{sec:7}.

锂离子电池的健康状态估计是电池管理系统的核心功能，电池的容量衰减与内阻增长直接影响电动汽车的续航里程\cite{a2020}。The state of health of a battery is difficult to measure; the data-driven approach uses the charging curve to estimate it, see Eq.~\eqref{eq:7} and Section~\ref{sec:7}.

锂离子电池的健康状态估计是电池管理系统的核心功能，电池的容量衰减与内阻增长直接影响电动汽车的续航里程\cite{a2020}。The state of health of a battery is difficult to measure; the data-driven approach uses the charging curve to estimate it, see Eq.~\eqref{eq:7} and Section~\ref{sec:7}.

\begin{equation}\label{eq:7}
  Q_t = Q_0 - \int_0^t \frac{I(\tau)}{3600}\,\mathrm{d}\tau
\end{equation}
\begin{table}[h]\centering\begin{tabular}{lcc}\toprule 方法 & RMSE & MAE\\\midrule GRU & 1.2 & 0.9\\\bottomrule\end{tabular}\caption{结果 7}\end{table}

\section{第 8 节 电池健康状态估计}\label{sec:8}
锂离子电池的健康状态估计是电池管理系统的核心功能，电池的容量衰减与内阻增长直接影响电动汽车的续航里程\cite{a2020}。The state of health of a battery is difficult to measure; the data-driven approach uses the charging curve to estimate it, see Eq.~\eqref{eq:8} and Section~\ref{sec:8}.

锂离子电池的健康状态估计是电池管理系统的核心功能，电池的容量衰减与内阻增长直接影响电动汽车的续航里程\cite{a2020}。The state of health of a battery is difficult to measure; the data-driven approach uses the charging curve to estimate it, see Eq.~\eqref{eq:8} and Section~\ref{sec:8}.

锂离子电池的健康状态估计是电池管理系统的核心功能，电池的容量衰减与内阻增长直接影响电动汽车的续航里程\cite{a2020}。The state of health of a battery is difficult to measure; the data-driven approach uses the charging curve to estimate it, see Eq.~\eqref{eq:8} and Section~\ref{sec:8}.

锂离子电池的健康状态估计是电池管理系统的核心功能，电池的容量衰减与内阻增长直接影响电动汽车的续航里程\cite{a2020}。The state of health of a battery is difficult to measure; the data-driven approach uses the charging curve to estimate it, see Eq.~\eqref{eq:8} and Section~\ref{sec:8}.

\begin{equation}\label{eq:8}
  Q_t = Q_0 - \int_0^t \frac{I(\tau)}{3600}\,\mathrm{d}\tau
\end{equation}
\begin{table}[h]\centering\begin{tabular}{lcc}\toprule 方法 & RMSE & MAE\\\midrule GRU & 1.2 & 0.9\\\bottomrule\end{tabular}\caption{结果 8}\end{table}

\section{第 9 节 电池健康状态估计}\label{sec:9}
锂离子电池的健康状态估计是电池管理系统的核心功能，电池的容量衰减与内阻增长直接影响电动汽车的续航里程\cite{a2020}。The state of health of a battery is difficult to measure; the data-driven approach uses the charging curve to estimate it, see Eq.~\eqref{eq:9} and Section~\ref{sec:9}.

锂离子电池的健康状态估计是电池管理系统的核心功能，电池的容量衰减与内阻增长直接影响电动汽车的续航里程\cite{a2020}。The state of health of a battery is difficult to measure; the data-driven approach uses the charging curve to estimate it, see Eq.~\eqref{eq:9} and Section~\ref{sec:9}.

锂离子电池的健康状态估计是电池管理系统的核心功能，电池的容量衰减与内阻增长直接影响电动汽车的续航里程\cite{a2020}。The state of health of a battery is difficult to measure; the data-driven approach uses the charging curve to estimate it, see Eq.~\eqref{eq:9} and Section~\ref{sec:9}.

锂离子电池的健康状态估计是电池管理系统的核心功能，电池的容量衰减与内阻增长直接影响电动汽车的续航里程\cite{a2020}。The state of health of a battery is difficult to measure; the data-driven approach uses the charging curve to estimate it, see Eq.~\eqref{eq:9} and Section~\ref{sec:9}.

\begin{equation}\label{eq:9}
  Q_t = Q_0 - \int_0^t \frac{I(\tau)}{3600}\,\mathrm{d}\tau
\end{equation}
\begin{table}[h]\centering\begin{tabular}{lcc}\toprule 方法 & RMSE & MAE\\\midrule GRU & 1.2 & 0.9\\\bottomrule\end{tabular}\caption{结果 9}\end{table}

\section{第 10 节 电池健康状态估计}\label{sec:10}
锂离子电池的健康状态估计是电池管理系统的核心功能，电池的容量衰减与内阻增长直接影响电动汽车的续航里程\cite{a2020}。The state of health of a battery is difficult to measure; the data-driven approach uses the charging curve to estimate it, see Eq.~\eqref{eq:10} and Section~\ref{sec:10}.

锂离子电池的健康状态估计是电池管理系统的核心功能，电池的容量衰减与内阻增长直接影响电动汽车的续航里程\cite{a2020}。The state of health of a battery is difficult to measure; the data-driven approach uses the charging curve to estimate it, see Eq.~\eqref{eq:10} and Section~\ref{sec:10}.

锂离子电池的健康状态估计是电池管理系统的核心功能，电池的容量衰减与内阻增长直接影响电动汽车的续航里程\cite{a2020}。The state of health of a battery is difficult to measure; the data-driven approach uses the charging curve to estimate it, see Eq.~\eqref{eq:10} and Section~\ref{sec:10}.

锂离子电池的健康状态估计是电池管理系统的核心功能，电池的容量衰减与内阻增长直接影响电动汽车的续航里程\cite{a2020}。The state of health of a battery is difficult to measure; the data-driven approach uses the charging curve to estimate it, see Eq.~\eqref{eq:10} and Section~\ref{sec:10}.

\begin{equation}\label{eq:10}
  Q_t = Q_0 - \int_0^t \frac{I(\tau)}{3600}\,\mathrm{d}\tau
\end{equation}
\begin{table}[h]\centering\begin{tabular}{lcc}\toprule 方法 & RMSE & MAE\\\midrule GRU & 1.2 & 0.9\\\bottomrule\end{tabular}\caption{结果 10}\end{table}

\section{第 11 节 电池健康状态估计}\label{sec:11}
锂离子电池的健康状态估计是电池管理系统的核心功能，电池的容量衰减与内阻增长直接影响电动汽车的续航里程\cite{a2020}。The state of health of a battery is difficult to measure; the data-driven approach uses the charging curve to estimate it, see Eq.~\eqref{eq:11} and Section~\ref{sec:11}.

锂离子电池的健康状态估计是电池管理系统的核心功能，电池的容量衰减与内阻增长直接影响电动汽车的续航里程\cite{a2020}。The state of health of a battery is difficult to measure; the data-driven approach uses the charging curve to estimate it, see Eq.~\eqref{eq:11} and Section~\ref{sec:11}.

锂离子电池的健康状态估计是电池管理系统的核心功能，电池的容量衰减与内阻增长直接影响电动汽车的续航里程\cite{a2020}。The state of health of a battery is difficult to measure; the data-driven approach uses the charging curve to estimate it, see Eq.~\eqref{eq:11} and Section~\ref{sec:11}.

锂离子电池的健康状态估计是电池管理系统的核心功能，电池的容量衰减与内阻增长直接影响电动汽车的续航里程\cite{a2020}。The state of health of a battery is difficult to measure; the data-driven approach uses the charging curve to estimate it, see Eq.~\eqref{eq:11} and Section~\ref{sec:11}.

\begin{equation}\label{eq:11}
  Q_t = Q_0 - \int_0^t \frac{I(\tau)}{3600}\,\mathrm{d}\tau
\end{equation}
\begin{table}[h]\centering\begin{tabular}{lcc}\toprule 方法 & RMSE & MAE\\\midrule GRU & 1.2 & 0.9\\\bottomrule\end{tabular}\caption{结果 11}\end{table}

\section{第 12 节 电池健康状态估计}\label{sec:12}
锂离子电池的健康状态估计是电池管理系统的核心功能，电池的容量衰减与内阻增长直接影响电动汽车的续航里程\cite{a2020}。The state of health of a battery is difficult to measure; the data-driven approach uses the charging curve to estimate it, see Eq.~\eqref{eq:12} and Section~\ref{sec:12}.

锂离子电池的健康状态估计是电池管理系统的核心功能，电池的容量衰减与内阻增长直接影响电动汽车的续航里程\cite{a2020}。The state of health of a battery is difficult to measure; the data-driven approach uses the charging curve to estimate it, see Eq.~\eqref{eq:12} and Section~\ref{sec:12}.

锂离子电池的健康状态估计是电池管理系统的核心功能，电池的容量衰减与内阻增长直接影响电动汽车的续航里程\cite{a2020}。The state of health of a battery is difficult to measure; the data-driven approach uses the charging curve to estimate it, see Eq.~\eqref{eq:12} and Section~\ref{sec:12}.

锂离子电池的健康状态估计是电池管理系统的核心功能，电池的容量衰减与内阻增长直接影响电动汽车的续航里程\cite{a2020}。The state of health of a battery is difficult to measure; the data-driven approach uses the charging curve to estimate it, see Eq.~\eqref{eq:12} and Section~\ref{sec:12}.

\begin{equation}\label{eq:12}
  Q_t = Q_0 - \int_0^t \frac{I(\tau)}{3600}\,\mathrm{d}\tau
\end{equation}
\begin{table}[h]\centering\begin{tabular}{lcc}\toprule 方法 & RMSE & MAE\\\midrule GRU & 1.2 & 0.9\\\bottomrule\end{tabular}\caption{结果 12}\end{table}

\section{第 13 节 电池健康状态估计}\label{sec:13}
锂离子电池的健康状态估计是电池管理系统的核心功能，电池的容量衰减与内阻增长直接影响电动汽车的续航里程\cite{a2020}。The state of health of a battery is difficult to measure; the data-driven approach uses the charging curve to estimate it, see Eq.~\eqref{eq:13} and Section~\ref{sec:13}.

锂离子电池的健康状态估计是电池管理系统的核心功能，电池的容量衰减与内阻增长直接影响电动汽车的续航里程\cite{a2020}。The state of health of a battery is difficult to measure; the data-driven approach uses the charging curve to estimate it, see Eq.~\eqref{eq:13} and Section~\ref{sec:13}.

锂离子电池的健康状态估计是电池管理系统的核心功能，电池的容量衰减与内阻增长直接影响电动汽车的续航里程\cite{a2020}。The state of health of a battery is difficult to measure; the data-driven approach uses the charging curve to estimate it, see Eq.~\eqref{eq:13} and Section~\ref{sec:13}.

锂离子电池的健康状态估计是电池管理系统的核心功能，电池的容量衰减与内阻增长直接影响电动汽车的续航里程\cite{a2020}。The state of health of a battery is difficult to measure; the data-driven approach uses the charging curve to estimate it, see Eq.~\eqref{eq:13} and Section~\ref{sec:13}.

\begin{equation}\label{eq:13}
  Q_t = Q_0 - \int_0^t \frac{I(\tau)}{3600}\,\mathrm{d}\tau
\end{equation}
\begin{table}[h]\centering\begin{tabular}{lcc}\toprule 方法 & RMSE & MAE\\\midrule GRU & 1.2 & 0.9\\\bottomrule\end{tabular}\caption{结果 13}\end{table}

\section{第 14 节 电池健康状态估计}\label{sec:14}
锂离子电池的健康状态估计是电池管理系统的核心功能，电池的容量衰减与内阻增长直接影响电动汽车的续航里程\cite{a2020}。The state of health of a battery is difficult to measure; the data-driven approach uses the charging curve to estimate it, see Eq.~\eqref{eq:14} and Section~\ref{sec:14}.

锂离子电池的健康状态估计是电池管理系统的核心功能，电池的容量衰减与内阻增长直接影响电动汽车的续航里程\cite{a2020}。The state of health of a battery is difficult to measure; the data-driven approach uses the charging curve to estimate it, see Eq.~\eqref{eq:14} and Section~\ref{sec:14}.

锂离子电池的健康状态估计是电池管理系统的核心功能，电池的容量衰减与内阻增长直接影响电动汽车的续航里程\cite{a2020}。The state of health of a battery is difficult to measure; the data-driven approach uses the charging curve to estimate it, see Eq.~\eqref{eq:14} and Section~\ref{sec:14}.

锂离子电池的健康状态估计是电池管理系统的核心功能，电池的容量衰减与内阻增长直接影响电动汽车的续航里程\cite{a2020}。The state of health of a battery is difficult to measure; the data-driven approach uses the charging curve to estimate it, see Eq.~\eqref{eq:14} and Section~\ref{sec:14}.

\begin{equation}\label{eq:14}
  Q_t = Q_0 - \int_0^t \frac{I(\tau)}{3600}\,\mathrm{d}\tau
\end{equation}
\begin{table}[h]\centering\begin{tabular}{lcc}\toprule 方法 & RMSE & MAE\\\midrule GRU & 1.2 & 0.9\\\bottomrule\end{tabular}\caption{结果 14}\end{table}

\section{第 15 节 电池健康状态估计}\label{sec:15}
锂离子电池的健康状态估计是电池管理系统的核心功能，电池的容量衰减与内阻增长直接影响电动汽车的续航里程\cite{a2020}。The state of health of a battery is difficult to measure; the data-driven approach uses the charging curve to estimate it, see Eq.~\eqref{eq:15} and Section~\ref{sec:15}.

锂离子电池的健康状态估计是电池管理系统的核心功能，电池的容量衰减与内阻增长直接影响电动汽车的续航里程\cite{a2020}。The state of health of a battery is difficult to measure; the data-driven approach uses the charging curve to estimate it, see Eq.~\eqref{eq:15} and Section~\ref{sec:15}.

锂离子电池的健康状态估计是电池管理系统的核心功能，电池的容量衰减与内阻增长直接影响电动汽车的续航里程\cite{a2020}。The state of health of a battery is difficult to measure; the data-driven approach uses the charging curve to estimate it, see Eq.~\eqref{eq:15} and Section~\ref{sec:15}.

锂离子电池的健康状态估计是电池管理系统的核心功能，电池的容量衰减与内阻增长直接影响电动汽车的续航里程\cite{a2020}。The state of health of a battery is difficult to measure; the data-driven approach uses the charging curve to estimate it, see Eq.~\eqref{eq:15} and Section~\ref{sec:15}.

\begin{equation}\label{eq:15}
  Q_t = Q_0 - \int_0^t \frac{I(\tau)}{3600}\,\mathrm{d}\tau
\end{equation}
\begin{table}[h]\centering\begin{tabular}{lcc}\toprule 方法 & RMSE & MAE\\\midrule GRU & 1.2 & 0.9\\\bottomrule\end{tabular}\caption{结果 15}\end{table}

\section{子文件}
这是子文件里的一段文字，用来检查跨文件的定位是否正确。

\begin{figure}[h]\centering\includegraphics[width=3cm]{figs/box.pdf}\caption{图}\end{figure}
The \soh{} again.
\section{第 16 节 电池健康状态估计}\label{sec:16}
锂离子电池的健康状态估计是电池管理系统的核心功能，电池的容量衰减与内阻增长直接影响电动汽车的续航里程\cite{a2020}。The state of health of a battery is difficult to measure; the data-driven approach uses the charging curve to estimate it, see Eq.~\eqref{eq:16} and Section~\ref{sec:16}.

锂离子电池的健康状态估计是电池管理系统的核心功能，电池的容量衰减与内阻增长直接影响电动汽车的续航里程\cite{a2020}。The state of health of a battery is difficult to measure; the data-driven approach uses the charging curve to estimate it, see Eq.~\eqref{eq:16} and Section~\ref{sec:16}.

锂离子电池的健康状态估计是电池管理系统的核心功能，电池的容量衰减与内阻增长直接影响电动汽车的续航里程\cite{a2020}。The state of health of a battery is difficult to measure; the data-driven approach uses the charging curve to estimate it, see Eq.~\eqref{eq:16} and Section~\ref{sec:16}.

锂离子电池的健康状态估计是电池管理系统的核心功能，电池的容量衰减与内阻增长直接影响电动汽车的续航里程\cite{a2020}。The state of health of a battery is difficult to measure; the data-driven approach uses the charging curve to estimate it, see Eq.~\eqref{eq:16} and Section~\ref{sec:16}.

\begin{equation}\label{eq:16}
  Q_t = Q_0 - \int_0^t \frac{I(\tau)}{3600}\,\mathrm{d}\tau
\end{equation}
\begin{table}[h]\centering\begin{tabular}{lcc}\toprule 方法 & RMSE & MAE\\\midrule GRU & 1.2 & 0.9\\\bottomrule\end{tabular}\caption{结果 16}\end{table}

\section{第 17 节 电池健康状态估计}\label{sec:17}
锂离子电池的健康状态估计是电池管理系统的核心功能，电池的容量衰减与内阻增长直接影响电动汽车的续航里程\cite{a2020}。The state of health of a battery is difficult to measure; the data-driven approach uses the charging curve to estimate it, see Eq.~\eqref{eq:17} and Section~\ref{sec:17}.

锂离子电池的健康状态估计是电池管理系统的核心功能，电池的容量衰减与内阻增长直接影响电动汽车的续航里程\cite{a2020}。The state of health of a battery is difficult to measure; the data-driven approach uses the charging curve to estimate it, see Eq.~\eqref{eq:17} and Section~\ref{sec:17}.

锂离子电池的健康状态估计是电池管理系统的核心功能，电池的容量衰减与内阻增长直接影响电动汽车的续航里程\cite{a2020}。The state of health of a battery is difficult to measure; the data-driven approach uses the charging curve to estimate it, see Eq.~\eqref{eq:17} and Section~\ref{sec:17}.

锂离子电池的健康状态估计是电池管理系统的核心功能，电池的容量衰减与内阻增长直接影响电动汽车的续航里程\cite{a2020}。The state of health of a battery is difficult to measure; the data-driven approach uses the charging curve to estimate it, see Eq.~\eqref{eq:17} and Section~\ref{sec:17}.

\begin{equation}\label{eq:17}
  Q_t = Q_0 - \int_0^t \frac{I(\tau)}{3600}\,\mathrm{d}\tau
\end{equation}
\begin{table}[h]\centering\begin{tabular}{lcc}\toprule 方法 & RMSE & MAE\\\midrule GRU & 1.2 & 0.9\\\bottomrule\end{tabular}\caption{结果 17}\end{table}

\section{第 18 节 电池健康状态估计}\label{sec:18}
锂离子电池的健康状态估计是电池管理系统的核心功能，电池的容量衰减与内阻增长直接影响电动汽车的续航里程\cite{a2020}。The state of health of a battery is difficult to measure; the data-driven approach uses the charging curve to estimate it, see Eq.~\eqref{eq:18} and Section~\ref{sec:18}.

锂离子电池的健康状态估计是电池管理系统的核心功能，电池的容量衰减与内阻增长直接影响电动汽车的续航里程\cite{a2020}。The state of health of a battery is difficult to measure; the data-driven approach uses the charging curve to estimate it, see Eq.~\eqref{eq:18} and Section~\ref{sec:18}.

锂离子电池的健康状态估计是电池管理系统的核心功能，电池的容量衰减与内阻增长直接影响电动汽车的续航里程\cite{a2020}。The state of health of a battery is difficult to measure; the data-driven approach uses the charging curve to estimate it, see Eq.~\eqref{eq:18} and Section~\ref{sec:18}.

锂离子电池的健康状态估计是电池管理系统的核心功能，电池的容量衰减与内阻增长直接影响电动汽车的续航里程\cite{a2020}。The state of health of a battery is difficult to measure; the data-driven approach uses the charging curve to estimate it, see Eq.~\eqref{eq:18} and Section~\ref{sec:18}.

\begin{equation}\label{eq:18}
  Q_t = Q_0 - \int_0^t \frac{I(\tau)}{3600}\,\mathrm{d}\tau
\end{equation}
\begin{table}[h]\centering\begin{tabular}{lcc}\toprule 方法 & RMSE & MAE\\\midrule GRU & 1.2 & 0.9\\\bottomrule\end{tabular}\caption{结果 18}\end{table}

\section{第 19 节 电池健康状态估计}\label{sec:19}
锂离子电池的健康状态估计是电池管理系统的核心功能，电池的容量衰减与内阻增长直接影响电动汽车的续航里程\cite{a2020}。The state of health of a battery is difficult to measure; the data-driven approach uses the charging curve to estimate it, see Eq.~\eqref{eq:19} and Section~\ref{sec:19}.

锂离子电池的健康状态估计是电池管理系统的核心功能，电池的容量衰减与内阻增长直接影响电动汽车的续航里程\cite{a2020}。The state of health of a battery is difficult to measure; the data-driven approach uses the charging curve to estimate it, see Eq.~\eqref{eq:19} and Section~\ref{sec:19}.

锂离子电池的健康状态估计是电池管理系统的核心功能，电池的容量衰减与内阻增长直接影响电动汽车的续航里程\cite{a2020}。The state of health of a battery is difficult to measure; the data-driven approach uses the charging curve to estimate it, see Eq.~\eqref{eq:19} and Section~\ref{sec:19}.

锂离子电池的健康状态估计是电池管理系统的核心功能，电池的容量衰减与内阻增长直接影响电动汽车的续航里程\cite{a2020}。The state of health of a battery is difficult to measure; the data-driven approach uses the charging curve to estimate it, see Eq.~\eqref{eq:19} and Section~\ref{sec:19}.

\begin{equation}\label{eq:19}
  Q_t = Q_0 - \int_0^t \frac{I(\tau)}{3600}\,\mathrm{d}\tau
\end{equation}
\begin{table}[h]\centering\begin{tabular}{lcc}\toprule 方法 & RMSE & MAE\\\midrule GRU & 1.2 & 0.9\\\bottomrule\end{tabular}\caption{结果 19}\end{table}

\section{第 20 节 电池健康状态估计}\label{sec:20}
锂离子电池的健康状态估计是电池管理系统的核心功能，电池的容量衰减与内阻增长直接影响电动汽车的续航里程\cite{a2020}。The state of health of a battery is difficult to measure; the data-driven approach uses the charging curve to estimate it, see Eq.~\eqref{eq:20} and Section~\ref{sec:20}.

锂离子电池的健康状态估计是电池管理系统的核心功能，电池的容量衰减与内阻增长直接影响电动汽车的续航里程\cite{a2020}。The state of health of a battery is difficult to measure; the data-driven approach uses the charging curve to estimate it, see Eq.~\eqref{eq:20} and Section~\ref{sec:20}.

锂离子电池的健康状态估计是电池管理系统的核心功能，电池的容量衰减与内阻增长直接影响电动汽车的续航里程\cite{a2020}。The state of health of a battery is difficult to measure; the data-driven approach uses the charging curve to estimate it, see Eq.~\eqref{eq:20} and Section~\ref{sec:20}.

锂离子电池的健康状态估计是电池管理系统的核心功能，电池的容量衰减与内阻增长直接影响电动汽车的续航里程\cite{a2020}。The state of health of a battery is difficult to measure; the data-driven approach uses the charging curve to estimate it, see Eq.~\eqref{eq:20} and Section~\ref{sec:20}.

\begin{equation}\label{eq:20}
  Q_t = Q_0 - \int_0^t \frac{I(\tau)}{3600}\,\mathrm{d}\tau
\end{equation}
\begin{table}[h]\centering\begin{tabular}{lcc}\toprule 方法 & RMSE & MAE\\\midrule GRU & 1.2 & 0.9\\\bottomrule\end{tabular}\caption{结果 20}\end{table}

\section{第 21 节 电池健康状态估计}\label{sec:21}
锂离子电池的健康状态估计是电池管理系统的核心功能，电池的容量衰减与内阻增长直接影响电动汽车的续航里程\cite{a2020}。The state of health of a battery is difficult to measure; the data-driven approach uses the charging curve to estimate it, see Eq.~\eqref{eq:21} and Section~\ref{sec:21}.

锂离子电池的健康状态估计是电池管理系统的核心功能，电池的容量衰减与内阻增长直接影响电动汽车的续航里程\cite{a2020}。The state of health of a battery is difficult to measure; the data-driven approach uses the charging curve to estimate it, see Eq.~\eqref{eq:21} and Section~\ref{sec:21}.

锂离子电池的健康状态估计是电池管理系统的核心功能，电池的容量衰减与内阻增长直接影响电动汽车的续航里程\cite{a2020}。The state of health of a battery is difficult to measure; the data-driven approach uses the charging curve to estimate it, see Eq.~\eqref{eq:21} and Section~\ref{sec:21}.

锂离子电池的健康状态估计是电池管理系统的核心功能，电池的容量衰减与内阻增长直接影响电动汽车的续航里程\cite{a2020}。The state of health of a battery is difficult to measure; the data-driven approach uses the charging curve to estimate it, see Eq.~\eqref{eq:21} and Section~\ref{sec:21}.

\begin{equation}\label{eq:21}
  Q_t = Q_0 - \int_0^t \frac{I(\tau)}{3600}\,\mathrm{d}\tau
\end{equation}
\begin{table}[h]\centering\begin{tabular}{lcc}\toprule 方法 & RMSE & MAE\\\midrule GRU & 1.2 & 0.9\\\bottomrule\end{tabular}\caption{结果 21}\end{table}

\section{第 22 节 电池健康状态估计}\label{sec:22}
锂离子电池的健康状态估计是电池管理系统的核心功能，电池的容量衰减与内阻增长直接影响电动汽车的续航里程\cite{a2020}。The state of health of a battery is difficult to measure; the data-driven approach uses the charging curve to estimate it, see Eq.~\eqref{eq:22} and Section~\ref{sec:22}.

锂离子电池的健康状态估计是电池管理系统的核心功能，电池的容量衰减与内阻增长直接影响电动汽车的续航里程\cite{a2020}。The state of health of a battery is difficult to measure; the data-driven approach uses the charging curve to estimate it, see Eq.~\eqref{eq:22} and Section~\ref{sec:22}.

锂离子电池的健康状态估计是电池管理系统的核心功能，电池的容量衰减与内阻增长直接影响电动汽车的续航里程\cite{a2020}。The state of health of a battery is difficult to measure; the data-driven approach uses the charging curve to estimate it, see Eq.~\eqref{eq:22} and Section~\ref{sec:22}.

锂离子电池的健康状态估计是电池管理系统的核心功能，电池的容量衰减与内阻增长直接影响电动汽车的续航里程\cite{a2020}。The state of health of a battery is difficult to measure; the data-driven approach uses the charging curve to estimate it, see Eq.~\eqref{eq:22} and Section~\ref{sec:22}.

\begin{equation}\label{eq:22}
  Q_t = Q_0 - \int_0^t \frac{I(\tau)}{3600}\,\mathrm{d}\tau
\end{equation}
\begin{table}[h]\centering\begin{tabular}{lcc}\toprule 方法 & RMSE & MAE\\\midrule GRU & 1.2 & 0.9\\\bottomrule\end{tabular}\caption{结果 22}\end{table}

\section{第 23 节 电池健康状态估计}\label{sec:23}
锂离子电池的健康状态估计是电池管理系统的核心功能，电池的容量衰减与内阻增长直接影响电动汽车的续航里程\cite{a2020}。The state of health of a battery is difficult to measure; the data-driven approach uses the charging curve to estimate it, see Eq.~\eqref{eq:23} and Section~\ref{sec:23}.

锂离子电池的健康状态估计是电池管理系统的核心功能，电池的容量衰减与内阻增长直接影响电动汽车的续航里程\cite{a2020}。The state of health of a battery is difficult to measure; the data-driven approach uses the charging curve to estimate it, see Eq.~\eqref{eq:23} and Section~\ref{sec:23}.

锂离子电池的健康状态估计是电池管理系统的核心功能，电池的容量衰减与内阻增长直接影响电动汽车的续航里程\cite{a2020}。The state of health of a battery is difficult to measure; the data-driven approach uses the charging curve to estimate it, see Eq.~\eqref{eq:23} and Section~\ref{sec:23}.

锂离子电池的健康状态估计是电池管理系统的核心功能，电池的容量衰减与内阻增长直接影响电动汽车的续航里程\cite{a2020}。The state of health of a battery is difficult to measure; the data-driven approach uses the charging curve to estimate it, see Eq.~\eqref{eq:23} and Section~\ref{sec:23}.

\begin{equation}\label{eq:23}
  Q_t = Q_0 - \int_0^t \frac{I(\tau)}{3600}\,\mathrm{d}\tau
\end{equation}
\begin{table}[h]\centering\begin{tabular}{lcc}\toprule 方法 & RMSE & MAE\\\midrule GRU & 1.2 & 0.9\\\bottomrule\end{tabular}\caption{结果 23}\end{table}

\section{第 24 节 电池健康状态估计}\label{sec:24}
锂离子电池的健康状态估计是电池管理系统的核心功能，电池的容量衰减与内阻增长直接影响电动汽车的续航里程\cite{a2020}。The state of health of a battery is difficult to measure; the data-driven approach uses the charging curve to estimate it, see Eq.~\eqref{eq:24} and Section~\ref{sec:24}.

锂离子电池的健康状态估计是电池管理系统的核心功能，电池的容量衰减与内阻增长直接影响电动汽车的续航里程\cite{a2020}。The state of health of a battery is difficult to measure; the data-driven approach uses the charging curve to estimate it, see Eq.~\eqref{eq:24} and Section~\ref{sec:24}.

锂离子电池的健康状态估计是电池管理系统的核心功能，电池的容量衰减与内阻增长直接影响电动汽车的续航里程\cite{a2020}。The state of health of a battery is difficult to measure; the data-driven approach uses the charging curve to estimate it, see Eq.~\eqref{eq:24} and Section~\ref{sec:24}.

锂离子电池的健康状态估计是电池管理系统的核心功能，电池的容量衰减与内阻增长直接影响电动汽车的续航里程\cite{a2020}。The state of health of a battery is difficult to measure; the data-driven approach uses the charging curve to estimate it, see Eq.~\eqref{eq:24} and Section~\ref{sec:24}.

\begin{equation}\label{eq:24}
  Q_t = Q_0 - \int_0^t \frac{I(\tau)}{3600}\,\mathrm{d}\tau
\end{equation}
\begin{table}[h]\centering\begin{tabular}{lcc}\toprule 方法 & RMSE & MAE\\\midrule GRU & 1.2 & 0.9\\\bottomrule\end{tabular}\caption{结果 24}\end{table}

\section{第 25 节 电池健康状态估计}\label{sec:25}
锂离子电池的健康状态估计是电池管理系统的核心功能，电池的容量衰减与内阻增长直接影响电动汽车的续航里程\cite{a2020}。The state of health of a battery is difficult to measure; the data-driven approach uses the charging curve to estimate it, see Eq.~\eqref{eq:25} and Section~\ref{sec:25}.

锂离子电池的健康状态估计是电池管理系统的核心功能，电池的容量衰减与内阻增长直接影响电动汽车的续航里程\cite{a2020}。The state of health of a battery is difficult to measure; the data-driven approach uses the charging curve to estimate it, see Eq.~\eqref{eq:25} and Section~\ref{sec:25}.

锂离子电池的健康状态估计是电池管理系统的核心功能，电池的容量衰减与内阻增长直接影响电动汽车的续航里程\cite{a2020}。The state of health of a battery is difficult to measure; the data-driven approach uses the charging curve to estimate it, see Eq.~\eqref{eq:25} and Section~\ref{sec:25}.

锂离子电池的健康状态估计是电池管理系统的核心功能，电池的容量衰减与内阻增长直接影响电动汽车的续航里程\cite{a2020}。The state of health of a battery is difficult to measure; the data-driven approach uses the charging curve to estimate it, see Eq.~\eqref{eq:25} and Section~\ref{sec:25}.

\begin{equation}\label{eq:25}
  Q_t = Q_0 - \int_0^t \frac{I(\tau)}{3600}\,\mathrm{d}\tau
\end{equation}
\begin{table}[h]\centering\begin{tabular}{lcc}\toprule 方法 & RMSE & MAE\\\midrule GRU & 1.2 & 0.9\\\bottomrule\end{tabular}\caption{结果 25}\end{table}

\section{第 26 节 电池健康状态估计}\label{sec:26}
锂离子电池的健康状态估计是电池管理系统的核心功能，电池的容量衰减与内阻增长直接影响电动汽车的续航里程\cite{a2020}。The state of health of a battery is difficult to measure; the data-driven approach uses the charging curve to estimate it, see Eq.~\eqref{eq:26} and Section~\ref{sec:26}.

锂离子电池的健康状态估计是电池管理系统的核心功能，电池的容量衰减与内阻增长直接影响电动汽车的续航里程\cite{a2020}。The state of health of a battery is difficult to measure; the data-driven approach uses the charging curve to estimate it, see Eq.~\eqref{eq:26} and Section~\ref{sec:26}.

锂离子电池的健康状态估计是电池管理系统的核心功能，电池的容量衰减与内阻增长直接影响电动汽车的续航里程\cite{a2020}。The state of health of a battery is difficult to measure; the data-driven approach uses the charging curve to estimate it, see Eq.~\eqref{eq:26} and Section~\ref{sec:26}.

锂离子电池的健康状态估计是电池管理系统的核心功能，电池的容量衰减与内阻增长直接影响电动汽车的续航里程\cite{a2020}。The state of health of a battery is difficult to measure; the data-driven approach uses the charging curve to estimate it, see Eq.~\eqref{eq:26} and Section~\ref{sec:26}.

\begin{equation}\label{eq:26}
  Q_t = Q_0 - \int_0^t \frac{I(\tau)}{3600}\,\mathrm{d}\tau
\end{equation}
\begin{table}[h]\centering\begin{tabular}{lcc}\toprule 方法 & RMSE & MAE\\\midrule GRU & 1.2 & 0.9\\\bottomrule\end{tabular}\caption{结果 26}\end{table}

\section{第 27 节 电池健康状态估计}\label{sec:27}
锂离子电池的健康状态估计是电池管理系统的核心功能，电池的容量衰减与内阻增长直接影响电动汽车的续航里程\cite{a2020}。The state of health of a battery is difficult to measure; the data-driven approach uses the charging curve to estimate it, see Eq.~\eqref{eq:27} and Section~\ref{sec:27}.

锂离子电池的健康状态估计是电池管理系统的核心功能，电池的容量衰减与内阻增长直接影响电动汽车的续航里程\cite{a2020}。The state of health of a battery is difficult to measure; the data-driven approach uses the charging curve to estimate it, see Eq.~\eqref{eq:27} and Section~\ref{sec:27}.

锂离子电池的健康状态估计是电池管理系统的核心功能，电池的容量衰减与内阻增长直接影响电动汽车的续航里程\cite{a2020}。The state of health of a battery is difficult to measure; the data-driven approach uses the charging curve to estimate it, see Eq.~\eqref{eq:27} and Section~\ref{sec:27}.

锂离子电池的健康状态估计是电池管理系统的核心功能，电池的容量衰减与内阻增长直接影响电动汽车的续航里程\cite{a2020}。The state of health of a battery is difficult to measure; the data-driven approach uses the charging curve to estimate it, see Eq.~\eqref{eq:27} and Section~\ref{sec:27}.

\begin{equation}\label{eq:27}
  Q_t = Q_0 - \int_0^t \frac{I(\tau)}{3600}\,\mathrm{d}\tau
\end{equation}
\begin{table}[h]\centering\begin{tabular}{lcc}\toprule 方法 & RMSE & MAE\\\midrule GRU & 1.2 & 0.9\\\bottomrule\end{tabular}\caption{结果 27}\end{table}

\section{第 28 节 电池健康状态估计}\label{sec:28}
锂离子电池的健康状态估计是电池管理系统的核心功能，电池的容量衰减与内阻增长直接影响电动汽车的续航里程\cite{a2020}。The state of health of a battery is difficult to measure; the data-driven approach uses the charging curve to estimate it, see Eq.~\eqref{eq:28} and Section~\ref{sec:28}.

锂离子电池的健康状态估计是电池管理系统的核心功能，电池的容量衰减与内阻增长直接影响电动汽车的续航里程\cite{a2020}。The state of health of a battery is difficult to measure; the data-driven approach uses the charging curve to estimate it, see Eq.~\eqref{eq:28} and Section~\ref{sec:28}.

锂离子电池的健康状态估计是电池管理系统的核心功能，电池的容量衰减与内阻增长直接影响电动汽车的续航里程\cite{a2020}。The state of health of a battery is difficult to measure; the data-driven approach uses the charging curve to estimate it, see Eq.~\eqref{eq:28} and Section~\ref{sec:28}.

锂离子电池的健康状态估计是电池管理系统的核心功能，电池的容量衰减与内阻增长直接影响电动汽车的续航里程\cite{a2020}。The state of health of a battery is difficult to measure; the data-driven approach uses the charging curve to estimate it, see Eq.~\eqref{eq:28} and Section~\ref{sec:28}.

\begin{equation}\label{eq:28}
  Q_t = Q_0 - \int_0^t \frac{I(\tau)}{3600}\,\mathrm{d}\tau
\end{equation}
\begin{table}[h]\centering\begin{tabular}{lcc}\toprule 方法 & RMSE & MAE\\\midrule GRU & 1.2 & 0.9\\\bottomrule\end{tabular}\caption{结果 28}\end{table}

\section{第 29 节 电池健康状态估计}\label{sec:29}
锂离子电池的健康状态估计是电池管理系统的核心功能，电池的容量衰减与内阻增长直接影响电动汽车的续航里程\cite{a2020}。The state of health of a battery is difficult to measure; the data-driven approach uses the charging curve to estimate it, see Eq.~\eqref{eq:29} and Section~\ref{sec:29}.

锂离子电池的健康状态估计是电池管理系统的核心功能，电池的容量衰减与内阻增长直接影响电动汽车的续航里程\cite{a2020}。The state of health of a battery is difficult to measure; the data-driven approach uses the charging curve to estimate it, see Eq.~\eqref{eq:29} and Section~\ref{sec:29}.

锂离子电池的健康状态估计是电池管理系统的核心功能，电池的容量衰减与内阻增长直接影响电动汽车的续航里程\cite{a2020}。The state of health of a battery is difficult to measure; the data-driven approach uses the charging curve to estimate it, see Eq.~\eqref{eq:29} and Section~\ref{sec:29}.

锂离子电池的健康状态估计是电池管理系统的核心功能，电池的容量衰减与内阻增长直接影响电动汽车的续航里程\cite{a2020}。The state of health of a battery is difficult to measure; the data-driven approach uses the charging curve to estimate it, see Eq.~\eqref{eq:29} and Section~\ref{sec:29}.

\begin{equation}\label{eq:29}
  Q_t = Q_0 - \int_0^t \frac{I(\tau)}{3600}\,\mathrm{d}\tau
\end{equation}
\begin{table}[h]\centering\begin{tabular}{lcc}\toprule 方法 & RMSE & MAE\\\midrule GRU & 1.2 & 0.9\\\bottomrule\end{tabular}\caption{结果 29}\end{table}

\section{第 30 节 电池健康状态估计}\label{sec:30}
锂离子电池的健康状态估计是电池管理系统的核心功能，电池的容量衰减与内阻增长直接影响电动汽车的续航里程\cite{a2020}。The state of health of a battery is difficult to measure; the data-driven approach uses the charging curve to estimate it, see Eq.~\eqref{eq:30} and Section~\ref{sec:30}.

锂离子电池的健康状态估计是电池管理系统的核心功能，电池的容量衰减与内阻增长直接影响电动汽车的续航里程\cite{a2020}。The state of health of a battery is difficult to measure; the data-driven approach uses the charging curve to estimate it, see Eq.~\eqref{eq:30} and Section~\ref{sec:30}.

锂离子电池的健康状态估计是电池管理系统的核心功能，电池的容量衰减与内阻增长直接影响电动汽车的续航里程\cite{a2020}。The state of health of a battery is difficult to measure; the data-driven approach uses the charging curve to estimate it, see Eq.~\eqref{eq:30} and Section~\ref{sec:30}.

锂离子电池的健康状态估计是电池管理系统的核心功能，电池的容量衰减与内阻增长直接影响电动汽车的续航里程\cite{a2020}。The state of health of a battery is difficult to measure; the data-driven approach uses the charging curve to estimate it, see Eq.~\eqref{eq:30} and Section~\ref{sec:30}.

\begin{equation}\label{eq:30}
  Q_t = Q_0 - \int_0^t \frac{I(\tau)}{3600}\,\mathrm{d}\tau
\end{equation}
\begin{table}[h]\centering\begin{tabular}{lcc}\toprule 方法 & RMSE & MAE\\\midrule GRU & 1.2 & 0.9\\\bottomrule\end{tabular}\caption{结果 30}\end{table}

\bibliographystyle{plain}
\bibliography{refs}
\end{document}
